\documentclass[conference]{IEEEtran}
\usepackage{graphicx,booktabs,amsmath,amssymb,cite,url}
\graphicspath{{../figures/}}
% AUTO-GENERATED by scripts/make_report_assets.py. Do not edit by hand.
\newcommand{\congOffpeak}{5.8}
\newcommand{\congPeak}{26.9}
\newcommand{\corrEightHop}{0.14}
\newcommand{\corrOneHop}{0.45}
\newcommand{\corrTwoHop}{0.25}
\newcommand{\gapAllFifteen}{0.059}
\newcommand{\gapAllFifteenHi}{0.069}
\newcommand{\gapAllFifteenLo}{0.051}
\newcommand{\gapAllSixty}{0.051}
\newcommand{\gapAllSixtyHi}{0.072}
\newcommand{\gapAllSixtyLo}{0.033}
\newcommand{\gapAllThirty}{0.073}
\newcommand{\gapAllThirtyHi}{0.090}
\newcommand{\gapAllThirtyLo}{0.059}
\newcommand{\gapCongestedFifteen}{0.336}
\newcommand{\gapCongestedFifteenHi}{0.392}
\newcommand{\gapCongestedFifteenLo}{0.282}
\newcommand{\gapCongestedSixty}{0.458}
\newcommand{\gapCongestedSixtyHi}{0.570}
\newcommand{\gapCongestedSixtyLo}{0.343}
\newcommand{\gapCongestedThirty}{0.559}
\newcommand{\gapCongestedThirtyHi}{0.670}
\newcommand{\gapCongestedThirtyLo}{0.452}
\newcommand{\gapFreeflowFifteen}{0.020}
\newcommand{\gapFreeflowFifteenHi}{0.024}
\newcommand{\gapFreeflowFifteenLo}{0.017}
\newcommand{\gapFreeflowSixty}{-0.007}
\newcommand{\gapFreeflowSixtyHi}{0.004}
\newcommand{\gapFreeflowSixtyLo}{-0.018}
\newcommand{\gapFreeflowThirty}{0.005}
\newcommand{\gapFreeflowThirtyHi}{0.012}
\newcommand{\gapFreeflowThirtyLo}{-0.003}
\newcommand{\gapOffpeakFifteen}{0.038}
\newcommand{\gapOffpeakFifteenHi}{0.048}
\newcommand{\gapOffpeakFifteenLo}{0.031}
\newcommand{\gapOffpeakSixty}{0.032}
\newcommand{\gapOffpeakSixtyHi}{0.055}
\newcommand{\gapOffpeakSixtyLo}{0.012}
\newcommand{\gapOffpeakThirty}{0.048}
\newcommand{\gapOffpeakThirtyHi}{0.065}
\newcommand{\gapOffpeakThirtyLo}{0.034}
\newcommand{\gapPeakFifteen}{0.117}
\newcommand{\gapPeakFifteenHi}{0.135}
\newcommand{\gapPeakFifteenLo}{0.101}
\newcommand{\gapPeakSixty}{0.112}
\newcommand{\gapPeakSixtyHi}{0.143}
\newcommand{\gapPeakSixtyLo}{0.081}
\newcommand{\gapPeakThirty}{0.152}
\newcommand{\gapPeakThirtyHi}{0.182}
\newcommand{\gapPeakThirtyLo}{0.126}
\newcommand{\gapWeekendFifteen}{0.036}
\newcommand{\gapWeekendFifteenHi}{0.057}
\newcommand{\gapWeekendFifteenLo}{0.025}
\newcommand{\gapWeekendSixty}{0.020}
\newcommand{\gapWeekendSixtyHi}{0.041}
\newcommand{\gapWeekendSixtyLo}{0.000}
\newcommand{\gapWeekendThirty}{0.037}
\newcommand{\gapWeekendThirtyHi}{0.057}
\newcommand{\gapWeekendThirtyLo}{0.024}
\newcommand{\graphDiameter}{17}
\newcommand{\graphEdges}{1515}
\newcommand{\graphNodes}{207}
\newcommand{\graphSCC}{195}
\newcommand{\maeHistoricalAverageFifteen}{4.19}
\newcommand{\maeHistoricalAverageFifteenCongested}{12.90}
\newcommand{\maeHistoricalAverageFifteenFreeflow}{2.96}
\newcommand{\maeHistoricalAverageFifteenOffpeak}{3.52}
\newcommand{\maeHistoricalAverageFifteenPeak}{6.04}
\newcommand{\maeHistoricalAverageFifteenPeakAm}{4.90}
\newcommand{\maeHistoricalAverageFifteenPeakPm}{7.13}
\newcommand{\maeHistoricalAverageFifteenWeekend}{3.50}
\newcommand{\maeHistoricalAverageSixty}{4.19}
\newcommand{\maeHistoricalAverageSixtyCongested}{12.90}
\newcommand{\maeHistoricalAverageSixtyFreeflow}{2.96}
\newcommand{\maeHistoricalAverageSixtyOffpeak}{3.51}
\newcommand{\maeHistoricalAverageSixtyPeak}{6.06}
\newcommand{\maeHistoricalAverageSixtyPeakAm}{4.93}
\newcommand{\maeHistoricalAverageSixtyPeakPm}{7.13}
\newcommand{\maeHistoricalAverageSixtyWeekend}{3.50}
\newcommand{\maeHistoricalAverageThirty}{4.19}
\newcommand{\maeHistoricalAverageThirtyCongested}{12.90}
\newcommand{\maeHistoricalAverageThirtyFreeflow}{2.96}
\newcommand{\maeHistoricalAverageThirtyOffpeak}{3.52}
\newcommand{\maeHistoricalAverageThirtyPeak}{6.05}
\newcommand{\maeHistoricalAverageThirtyPeakAm}{4.91}
\newcommand{\maeHistoricalAverageThirtyPeakPm}{7.13}
\newcommand{\maeHistoricalAverageThirtyWeekend}{3.50}
\newcommand{\maePersistenceFifteen}{3.49}
\newcommand{\maePersistenceFifteenCongested}{7.31}
\newcommand{\maePersistenceFifteenFreeflow}{2.96}
\newcommand{\maePersistenceFifteenOffpeak}{3.29}
\newcommand{\maePersistenceFifteenPeak}{4.39}
\newcommand{\maePersistenceFifteenPeakAm}{3.75}
\newcommand{\maePersistenceFifteenPeakPm}{4.99}
\newcommand{\maePersistenceFifteenWeekend}{2.86}
\newcommand{\maePersistenceSixty}{5.37}
\newcommand{\maePersistenceSixtyCongested}{13.33}
\newcommand{\maePersistenceSixtyFreeflow}{4.24}
\newcommand{\maePersistenceSixtyOffpeak}{4.85}
\newcommand{\maePersistenceSixtyPeak}{7.45}
\newcommand{\maePersistenceSixtyPeakAm}{6.93}
\newcommand{\maePersistenceSixtyPeakPm}{7.95}
\newcommand{\maePersistenceSixtyWeekend}{3.98}
\newcommand{\maePersistenceThirty}{4.20}
\newcommand{\maePersistenceThirtyCongested}{9.63}
\newcommand{\maePersistenceThirtyFreeflow}{3.43}
\newcommand{\maePersistenceThirtyOffpeak}{3.84}
\newcommand{\maePersistenceThirtyPeak}{5.58}
\newcommand{\maePersistenceThirtyPeakAm}{4.96}
\newcommand{\maePersistenceThirtyPeakPm}{6.17}
\newcommand{\maePersistenceThirtyWeekend}{3.31}
\newcommand{\maeRfMultioutputFifteen}{3.14}
\newcommand{\maeRfMultioutputFifteenCongested}{7.68}
\newcommand{\maeRfMultioutputFifteenFreeflow}{2.49}
\newcommand{\maeRfMultioutputFifteenOffpeak}{2.80}
\newcommand{\maeRfMultioutputFifteenPeak}{4.21}
\newcommand{\maeRfMultioutputFifteenPeakAm}{3.59}
\newcommand{\maeRfMultioutputFifteenPeakPm}{4.81}
\newcommand{\maeRfMultioutputFifteenWeekend}{2.59}
\newcommand{\maeRfMultioutputSixty}{4.73}
\newcommand{\maeRfMultioutputSixtyCongested}{14.05}
\newcommand{\maeRfMultioutputSixtyFreeflow}{3.42}
\newcommand{\maeRfMultioutputSixtyOffpeak}{3.82}
\newcommand{\maeRfMultioutputSixtyPeak}{7.24}
\newcommand{\maeRfMultioutputSixtyPeakAm}{6.52}
\newcommand{\maeRfMultioutputSixtyPeakPm}{7.91}
\newcommand{\maeRfMultioutputSixtyWeekend}{3.80}
\newcommand{\maeRfMultioutputThirty}{3.80}
\newcommand{\maeRfMultioutputThirtyCongested}{10.29}
\newcommand{\maeRfMultioutputThirtyFreeflow}{2.88}
\newcommand{\maeRfMultioutputThirtyOffpeak}{3.24}
\newcommand{\maeRfMultioutputThirtyPeak}{5.44}
\newcommand{\maeRfMultioutputThirtyPeakAm}{4.79}
\newcommand{\maeRfMultioutputThirtyPeakPm}{6.05}
\newcommand{\maeRfMultioutputThirtyWeekend}{3.08}
\newcommand{\maeSeasonalNaiveFifteen}{5.50}
\newcommand{\maeSeasonalNaiveFifteenCongested}{14.79}
\newcommand{\maeSeasonalNaiveFifteenFreeflow}{4.19}
\newcommand{\maeSeasonalNaiveFifteenOffpeak}{4.38}
\newcommand{\maeSeasonalNaiveFifteenPeak}{7.45}
\newcommand{\maeSeasonalNaiveFifteenPeakAm}{6.14}
\newcommand{\maeSeasonalNaiveFifteenPeakPm}{8.70}
\newcommand{\maeSeasonalNaiveFifteenWeekend}{5.83}
\newcommand{\maeSeasonalNaiveSixty}{5.50}
\newcommand{\maeSeasonalNaiveSixtyCongested}{14.80}
\newcommand{\maeSeasonalNaiveSixtyFreeflow}{4.19}
\newcommand{\maeSeasonalNaiveSixtyOffpeak}{4.37}
\newcommand{\maeSeasonalNaiveSixtyPeak}{7.48}
\newcommand{\maeSeasonalNaiveSixtyPeakAm}{6.18}
\newcommand{\maeSeasonalNaiveSixtyPeakPm}{8.70}
\newcommand{\maeSeasonalNaiveSixtyWeekend}{5.83}
\newcommand{\maeSeasonalNaiveThirty}{5.50}
\newcommand{\maeSeasonalNaiveThirtyCongested}{14.79}
\newcommand{\maeSeasonalNaiveThirtyFreeflow}{4.19}
\newcommand{\maeSeasonalNaiveThirtyOffpeak}{4.38}
\newcommand{\maeSeasonalNaiveThirtyPeak}{7.46}
\newcommand{\maeSeasonalNaiveThirtyPeakAm}{6.15}
\newcommand{\maeSeasonalNaiveThirtyPeakPm}{8.70}
\newcommand{\maeSeasonalNaiveThirtyWeekend}{5.83}
\newcommand{\maeXgbGlobalFifteen}{2.88}
\newcommand{\maeXgbGlobalFifteenCongested}{7.51}
\newcommand{\maeXgbGlobalFifteenFreeflow}{2.23}
\newcommand{\maeXgbGlobalFifteenOffpeak}{2.58}
\newcommand{\maeXgbGlobalFifteenPeak}{3.87}
\newcommand{\maeXgbGlobalFifteenPeakAm}{3.21}
\newcommand{\maeXgbGlobalFifteenPeakPm}{4.49}
\newcommand{\maeXgbGlobalFifteenWeekend}{2.36}
\newcommand{\maeXgbGlobalSixty}{3.78}
\newcommand{\maeXgbGlobalSixtyCongested}{13.03}
\newcommand{\maeXgbGlobalSixtyFreeflow}{2.48}
\newcommand{\maeXgbGlobalSixtyOffpeak}{3.11}
\newcommand{\maeXgbGlobalSixtyPeak}{5.65}
\newcommand{\maeXgbGlobalSixtyPeakAm}{4.67}
\newcommand{\maeXgbGlobalSixtyPeakPm}{6.58}
\newcommand{\maeXgbGlobalSixtyWeekend}{3.05}
\newcommand{\maeXgbGlobalThirty}{3.33}
\newcommand{\maeXgbGlobalThirtyCongested}{10.15}
\newcommand{\maeXgbGlobalThirtyFreeflow}{2.37}
\newcommand{\maeXgbGlobalThirtyOffpeak}{2.86}
\newcommand{\maeXgbGlobalThirtyPeak}{4.73}
\newcommand{\maeXgbGlobalThirtyPeakAm}{3.99}
\newcommand{\maeXgbGlobalThirtyPeakPm}{5.44}
\newcommand{\maeXgbGlobalThirtyWeekend}{2.70}
\newcommand{\maeXgbPerSegmentFifteen}{2.86}
\newcommand{\maeXgbPerSegmentFifteenCongested}{7.63}
\newcommand{\maeXgbPerSegmentFifteenFreeflow}{2.19}
\newcommand{\maeXgbPerSegmentFifteenOffpeak}{2.57}
\newcommand{\maeXgbPerSegmentFifteenPeak}{3.78}
\newcommand{\maeXgbPerSegmentFifteenPeakAm}{3.06}
\newcommand{\maeXgbPerSegmentFifteenPeakPm}{4.48}
\newcommand{\maeXgbPerSegmentFifteenWeekend}{2.38}
\newcommand{\maeXgbPerSegmentSixty}{3.68}
\newcommand{\maeXgbPerSegmentSixtyCongested}{12.15}
\newcommand{\maeXgbPerSegmentSixtyFreeflow}{2.49}
\newcommand{\maeXgbPerSegmentSixtyOffpeak}{3.11}
\newcommand{\maeXgbPerSegmentSixtyPeak}{5.26}
\newcommand{\maeXgbPerSegmentSixtyPeakAm}{4.19}
\newcommand{\maeXgbPerSegmentSixtyPeakPm}{6.28}
\newcommand{\maeXgbPerSegmentSixtyWeekend}{3.09}
\newcommand{\maeXgbPerSegmentThirty}{3.28}
\newcommand{\maeXgbPerSegmentThirtyCongested}{9.91}
\newcommand{\maeXgbPerSegmentThirtyFreeflow}{2.34}
\newcommand{\maeXgbPerSegmentThirtyOffpeak}{2.86}
\newcommand{\maeXgbPerSegmentThirtyPeak}{4.51}
\newcommand{\maeXgbPerSegmentThirtyPeakAm}{3.63}
\newcommand{\maeXgbPerSegmentThirtyPeakPm}{5.35}
\newcommand{\maeXgbPerSegmentThirtyWeekend}{2.73}
\newcommand{\maeXgbSpatialFifteen}{2.82}
\newcommand{\maeXgbSpatialFifteenCongested}{7.18}
\newcommand{\maeXgbSpatialFifteenFreeflow}{2.21}
\newcommand{\maeXgbSpatialFifteenOffpeak}{2.54}
\newcommand{\maeXgbSpatialFifteenPeak}{3.75}
\newcommand{\maeXgbSpatialFifteenPeakAm}{3.14}
\newcommand{\maeXgbSpatialFifteenPeakPm}{4.33}
\newcommand{\maeXgbSpatialFifteenWeekend}{2.32}
\newcommand{\maeXgbSpatialSixty}{3.73}
\newcommand{\maeXgbSpatialSixtyCongested}{12.58}
\newcommand{\maeXgbSpatialSixtyFreeflow}{2.48}
\newcommand{\maeXgbSpatialSixtyOffpeak}{3.08}
\newcommand{\maeXgbSpatialSixtyPeak}{5.54}
\newcommand{\maeXgbSpatialSixtyPeakAm}{4.65}
\newcommand{\maeXgbSpatialSixtyPeakPm}{6.38}
\newcommand{\maeXgbSpatialSixtyWeekend}{3.03}
\newcommand{\maeXgbSpatialThirty}{3.26}
\newcommand{\maeXgbSpatialThirtyCongested}{9.59}
\newcommand{\maeXgbSpatialThirtyFreeflow}{2.36}
\newcommand{\maeXgbSpatialThirtyOffpeak}{2.82}
\newcommand{\maeXgbSpatialThirtyPeak}{4.58}
\newcommand{\maeXgbSpatialThirtyPeakAm}{3.88}
\newcommand{\maeXgbSpatialThirtyPeakPm}{5.24}
\newcommand{\maeXgbSpatialThirtyWeekend}{2.66}
\newcommand{\mapeHistoricalAverageFifteen}{13.1}
\newcommand{\mapeHistoricalAverageSixty}{13.1}
\newcommand{\mapeHistoricalAverageThirty}{13.1}
\newcommand{\mapePersistenceFifteen}{8.6}
\newcommand{\mapePersistenceSixty}{14.5}
\newcommand{\mapePersistenceThirty}{10.8}
\newcommand{\mapeRfMultioutputFifteen}{8.5}
\newcommand{\mapeRfMultioutputSixty}{14.3}
\newcommand{\mapeRfMultioutputThirty}{10.9}
\newcommand{\mapeSeasonalNaiveFifteen}{15.6}
\newcommand{\mapeSeasonalNaiveSixty}{15.6}
\newcommand{\mapeSeasonalNaiveThirty}{15.6}
\newcommand{\mapeXgbGlobalFifteen}{7.9}
\newcommand{\mapeXgbGlobalSixty}{12.1}
\newcommand{\mapeXgbGlobalThirty}{9.9}
\newcommand{\mapeXgbPerSegmentFifteen}{7.9}
\newcommand{\mapeXgbPerSegmentSixty}{11.7}
\newcommand{\mapeXgbPerSegmentThirty}{9.8}
\newcommand{\mapeXgbSpatialFifteen}{7.6}
\newcommand{\mapeXgbSpatialSixty}{11.8}
\newcommand{\mapeXgbSpatialThirty}{9.5}
\newcommand{\meanSpeed}{58.5}
\newcommand{\missingRate}{8.1}
\newcommand{\missingTest}{12.1}
\newcommand{\missingTrain}{7.1}
\newcommand{\missingVal}{7.0}
\newcommand{\networkOutages}{2148}
\newcommand{\rmseHistoricalAverageFifteen}{7.86}
\newcommand{\rmseHistoricalAverageSixty}{7.86}
\newcommand{\rmseHistoricalAverageThirty}{7.86}
\newcommand{\rmsePersistenceFifteen}{6.39}
\newcommand{\rmsePersistenceSixty}{10.30}
\newcommand{\rmsePersistenceThirty}{8.02}
\newcommand{\rmseRfMultioutputFifteen}{5.70}
\newcommand{\rmseRfMultioutputSixty}{8.25}
\newcommand{\rmseRfMultioutputThirty}{6.89}
\newcommand{\rmseSeasonalNaiveFifteen}{10.71}
\newcommand{\rmseSeasonalNaiveSixty}{10.71}
\newcommand{\rmseSeasonalNaiveThirty}{10.71}
\newcommand{\rmseXgbGlobalFifteen}{5.69}
\newcommand{\rmseXgbGlobalSixty}{7.76}
\newcommand{\rmseXgbGlobalThirty}{6.80}
\newcommand{\rmseXgbPerSegmentFifteen}{5.60}
\newcommand{\rmseXgbPerSegmentSixty}{7.52}
\newcommand{\rmseXgbPerSegmentThirty}{6.62}
\newcommand{\rmseXgbSpatialFifteen}{5.46}
\newcommand{\rmseXgbSpatialSixty}{7.54}
\newcommand{\rmseXgbSpatialThirty}{6.54}
\newcommand{\splitTestEnd}{2012-06-27}
\newcommand{\splitTestStart}{2012-06-04}
\newcommand{\splitTrainEnd}{2012-05-23}
\newcommand{\splitTrainStart}{2012-03-01}
\newcommand{\splitValEnd}{2012-06-04}
\newcommand{\splitValStart}{2012-05-23}
\newcommand{\xcorrAllNeg}{30}
\newcommand{\xcorrAllPos}{38}
\newcommand{\xcorrAllZero}{32}
\newcommand{\xcorrPairs}{1228}
\newcommand{\xcorrPeakNeg}{38}
\newcommand{\xcorrPeakPos}{50}
\newcommand{\xcorrPeakZero}{12}

\begin{document}
\title{Graph Attention Networks and a Leakage-Safe Hybrid for Traffic Speed Forecasting on METR-LA (Phase 2)}
\author{\IEEEauthorblockN{Pranav Singh, Kunal, Himanshu, Kanchan, and Vaibhav}}
\maketitle

\begin{abstract}
Phase 1 showed that segment-level machine learning is a strong baseline and that hand-made neighbour features help modestly.
Phase 2 asks whether a learned spatio-temporal graph attention network (STGAT), its no-graph counterpart, and a hybrid that corrects XGBoost with a graph residual do better, and where.
On a strict chronological protocol with masked metrics, STGAT beats persistence at every horizon (\maeStgatFifteen{}/\maeStgatThirty{}/\maeStgatSixty{} mph vs.\ \maePersistenceFifteen{}/\maePersistenceThirty{}/\maePersistenceSixty{} at 15/30/60 min) but is on par with the no-graph GRU at short horizons; its edge appears at 60 minutes (\maeStgatSixty{} vs.\ \maeGruNographSixty{} mph).
The best models are hybrids: a residual hybrid at 15 and 30 minutes (\maeHybridResidualFifteen{}, \maeHybridResidualThirty{} mph) and a learned gate at 60 minutes (\maeHybridGatingSixty{} mph), with gains concentrated in congested states.
A simulated ramp-metering study (offline, synthetic network) shows reduced mainline congestion but no change in total time spent.
\end{abstract}

\section{Introduction}
Congestion spreads along the road network, so graph models should help most at long horizons and in peaks.
We test this with controlled comparisons: graph vs.\ no-graph deep models, graph attention vs.\ fixed weights, and a hybrid designed to learn only what the per-segment model misses.
All numbers are generated by scripts; failures and null results are reported as such.

\section{Deep Learning Literature}
DCRNN \cite{li2018dcrnn} couples diffusion convolution on the directed road graph with a GRU encoder--decoder; STGCN \cite{yu2018stgcn} uses gated temporal and graph convolutions; Graph WaveNet \cite{wu2019graphwavenet} learns an adaptive adjacency; ASTGCN \cite{guo2019astgcn} and GMAN \cite{zheng2020gman} add attention.
GCN \cite{kipf2017gcn} aggregates with fixed normalised weights; GAT \cite{velickovic2018gat} learns neighbour weights, and GATv2 \cite{brody2022gatv2} makes the attention dynamic.

\section{Data, Graph and Regularisation}
Data and splits follow Phase 1 (\splitTrainStart--\splitTestEnd{}; strict chronological split, masked losses, train-only scaling and imputation).
The graph is rebuilt from road distances with a thresholded Gaussian kernel and matches the published matrix; it has \graphEdges{} directed edges.
The default model uses the bidirectional graph $W+W^\top$; forward-only and backward-only graphs are ablations.
Regularisation: dropout (0.1), attention dropout (0.1), weight decay ($10^{-4}$), early stopping on validation MAE, gradient clipping; each is toggled in an ablation.

\section{Architecture}
\textbf{GAT layer.} For a node $i$ with neighbours $\mathcal N(i)\cup\{i\}$: $e_{ij}=\mathrm{LeakyReLU}(a^\top[Wh_i\Vert Wh_j])$, $\alpha_{ij}=\mathrm{softmax}_j(e_{ij})$, $h_i'=\sum_j\alpha_{ij}Wh_j$, with $K=4$ heads concatenated.
We implement it ourselves (dense masked form and an equivalent neighbour-list form) and test numerical equality with PyTorch Geometric's \texttt{GATConv} and \texttt{GATv2Conv}.
\textbf{STGAT.} Linear embedding, learnable node embedding, two blocks of (gated causal temporal convolution $\to$ multi-head GAT $\to$ ELU/dropout $\to$ residual $+$ LayerNorm), last-step readout and an MLP that outputs all 12 horizons at once.
\textbf{GRU-nograph.} The same embedding and decoder with a GRU per sensor and no spatial layer.
\textbf{Compute-driven simplification.} To run on a laptop, the graph layers act on the last 4 of 12 time steps and each epoch uses a random third of the windows; this applies equally to the GRU baseline's schedule.

\section{Hybrid Method}
XGBoost handles local autoregression and calendar effects; STGAT is trained on the residual $y-\hat y_{\mathrm{XGB}}$ and receives $\hat y_{\mathrm{XGB}}$ as a decoder input, with a zero-initialised last layer so that training starts at the XGBoost forecast.
Residuals must come from \emph{out-of-fold} forecasts: XGBoost's in-sample residuals are too small.
We split training time into five blocks and predict block $k\ge2$ from a model trained on earlier blocks, with a gap of $H+F-1$ anchors; block 1 is dropped.
Alternatives: a gate $\hat y=w\hat y_{\mathrm{STGAT}}+(1-w)\hat y_{\mathrm{XGB}}$ trained on the first 80\% of validation and stopped on the last 20\%, and the plain average.

\section{Experimental Setup}
Same split, masks and history ($H=12$) for every model; STGAT, GRU, and hybrids use three seeds (mean $\pm$ std); at most 40 epochs with patience 8; Adam, learning rate $10^{-3}$.
Significance: Diebold--Mariano with a HAC variance \cite{diebold1995comparing,newey1987simple,harvey1997testing} per seed pair (we report the statistic's median and the largest $p$ across seeds in the code outputs).
Limited compute means ML tuning budgets were small (15 trials) and pooled XGBoost used 25\% of the rows.

\section{Results}
Table~\ref{tab:main} and Fig.~\ref{fig:all} summarise test errors; Table~\ref{tab:sig} gives the DM statistics.
\begin{table}[t]
\centering
\caption{Test-set error (mph) by model and horizon (minutes).}\label{tab:main}
\footnotesize\resizebox{\columnwidth}{!}{%
\begin{tabular}{lrrrrrr}
\toprule
model & MAE@15 & MAE@30 & MAE@60 & RMSE@15 & RMSE@30 & RMSE@60 \\
\midrule
historical average & 4.19 & 4.19 & 4.19 & 7.86 & 7.86 & 7.86 \\
persistence & 3.49 & 4.20 & 5.37 & 6.39 & 8.02 & 10.30 \\
rf multioutput & 3.14 & 3.80 & 4.73 & 5.70 & 6.89 & 8.25 \\
seasonal naive & 5.50 & 5.50 & 5.50 & 10.71 & 10.71 & 10.71 \\
xgb global & 2.88 & 3.33 & 3.78 & 5.69 & 6.80 & 7.76 \\
xgb per segment & 2.86 & 3.28 & \textbf{3.68} & 5.60 & 6.62 & \textbf{7.52} \\
xgb spatial & \textbf{2.82} & \textbf{3.26} & 3.73 & \textbf{5.46} & \textbf{6.54} & 7.54 \\
\bottomrule
\end{tabular}}
\end{table}

\begin{figure}[t]\centering
\includegraphics[width=\columnwidth]{error_vs_horizon_all_models.pdf}
\caption{Test MAE versus horizon for all models (bands: $\pm$ std over seeds).}\label{fig:all}
\end{figure}
\textbf{Graph vs.\ no graph.} STGAT and the GRU are indistinguishable at 15 minutes (\maeStgatFifteen{} vs.\ \maeGruNographFifteen{} mph); STGAT is slightly ahead at 30 (\maeStgatThirty{} vs.\ \maeGruNographThirty{}, not significant) and clearly at 60 minutes (\maeStgatSixty{} vs.\ \maeGruNographSixty{} mph, DM significant).
\textbf{Deep vs.\ trees.} STGAT is slightly worse than pooled XGBoost at 15 minutes (\maeStgatFifteen{} vs.\ \maeXgbGlobalFifteen{}) and better at 60 (\maeStgatSixty{} vs.\ \maeXgbGlobalSixty{}); the per-segment XGBoost model remains very competitive (\maeXgbPerSegmentSixty{} at 60 minutes).
\textbf{Hybrids.} Every hybrid beats XGBoost-global at all three horizons. The residual hybrid is best at 15 and 30 minutes; the gate is best at 60 minutes (\maeHybridGatingSixty{} mph vs.\ \maeStgatSixty{} for STGAT). The residual hybrid's gain over XGBoost-global grows with horizon (\gainXgbGlobalAllFifteen{}, \gainXgbGlobalAllThirty{}, \gainXgbGlobalAllSixty{} mph) and is largest in congestion at 60 minutes (\gainXgbGlobalCongestedSixty{} mph; \gainXgbGlobalPeakSixty{} at peak; \gainXgbGlobalFreeflowSixty{} in free flow), where it slightly hurts.
Over STGAT its gain shrinks with horizon (\gainStgatAllFifteen{}, \gainStgatAllThirty{}, \gainStgatAllSixty{} mph), consistent with STGAT learning more of the local structure at long range.
The learned gate puts more weight on STGAT at peaks than off-peak (at 60 min: \gatePeakSixty{} vs.\ \gateOffpeakSixty{}), as hypothesised.

\begin{table}[t]
\centering
\caption{Diebold-Mariano statistics (negative: first model better; * p<0.05).}\label{tab:sig}
\footnotesize\resizebox{\ifdim\width>\columnwidth\columnwidth\else\width\fi}{!}{%
\begin{tabular}{lrrr}
\toprule
pair & DM@15 & DM@30 & DM@60 \\
\midrule
hybrid average vs stgat & -25.5* & -6.0* & -3.1* \\
hybrid average vs xgb global & -10.9* & -11.8* & -10.9* \\
hybrid gating vs stgat & -25.2* & -13.0* & -11.3* \\
hybrid gating vs xgb global & -12.6* & -10.6* & -9.9* \\
hybrid residual vs hybrid average & -7.2* & -5.2* & 1.1 \\
hybrid residual vs stgat & -21.3* & -9.4* & -1.3 \\
hybrid residual vs xgb global & -12.3* & -10.6* & -9.2* \\
stgat vs gru nograph & 0.4 & -6.2 & -6.8* \\
stgat vs xgb global & 4.1* & -4.7* & -5.5* \\
stgat vs xgb per segment & 6.5* & 0.3 & -0.7 \\
stgat vs xgb spatial & 17.7* & 2.4 & -3.7* \\
xgb global vs persistence & -50.4* & -34.8* & -23.0* \\
xgb global vs xgb per segment & 3.8* & 5.3* & 6.3* \\
xgb spatial vs xgb global & -14.4* & -11.5* & -6.9* \\
xgb spatial vs xgb per segment & -6.3* & -1.7 & 3.0* \\
\bottomrule
\end{tabular}}
\end{table}


\begin{figure}[t]\centering
\includegraphics[width=\columnwidth]{peak_offpeak_bars.pdf}
\caption{Error by slice at 15/30/60 minutes.}\label{fig:peak}
\end{figure}

\section{Ablations}
Each ablation changes one setting from the default and was run with one seed and a reduced epoch budget (the default was re-run with the same budget as reference).
\begin{table}[t]
\centering
\caption{Ablations (1 seed, 12 epochs): MAE (mph) and change vs. default.}\label{tab:abl}
\footnotesize\resizebox{\columnwidth}{!}{%
\begin{tabular}{lrrrrrr}
\toprule
setting & MAE@15 & MAE@30 & MAE@60 & d@15 & d@30 & d@60 \\
\midrule
A0 default & 2.98 & 3.41 & 3.95 & 0.00 & 0.00 & 0.00 \\
A1 identity & 2.95 & 3.37 & 3.82 & -0.03 & -0.03 & -0.14 \\
A2 random & 3.03 & 3.48 & 4.04 & 0.04 & 0.07 & 0.09 \\
A3 backward & 2.99 & 3.42 & 3.96 & 0.01 & 0.01 & 0.00 \\
A3 forward & 3.00 & 3.42 & 3.94 & 0.01 & 0.01 & -0.01 \\
A4 gatv2 & 2.95 & 3.38 & 3.93 & -0.03 & -0.03 & -0.02 \\
A4 mean & 2.96 & 3.38 & 3.89 & -0.03 & -0.02 & -0.06 \\
A5 heads1 & 2.97 & 3.39 & 3.90 & -0.02 & -0.02 & -0.06 \\
A5 heads2 & 2.96 & 3.38 & 3.92 & -0.02 & -0.03 & -0.04 \\
A5 heads8 & 2.97 & 3.40 & 3.95 & -0.01 & -0.00 & -0.01 \\
A6 layers1 & 2.99 & 3.44 & 4.02 & 0.01 & 0.03 & 0.07 \\
A6 layers3 & 2.98 & 3.41 & 3.96 & -0.01 & 0.00 & 0.01 \\
A7 no node emb & 3.13 & 3.73 & 4.58 & 0.15 & 0.32 & 0.63 \\
A8 no time & 3.08 & 3.69 & 4.64 & 0.09 & 0.28 & 0.69 \\
A9 zero mask & 3.00 & 3.42 & 3.96 & 0.02 & 0.02 & 0.01 \\
A10 no dropout & 2.95 & 3.37 & 3.87 & -0.04 & -0.04 & -0.09 \\
A10 no wd & 2.93 & 3.31 & 3.73 & -0.06 & -0.10 & -0.22 \\
A11 H24 & 2.98 & 3.40 & 3.89 & -0.01 & -0.01 & -0.06 \\
A11 H6 & 2.98 & 3.40 & 3.93 & -0.01 & -0.01 & -0.02 \\
A12 edge bias & 2.95 & 3.37 & 3.90 & -0.03 & -0.04 & -0.06 \\
\bottomrule
\end{tabular}}
\end{table}

Under this reduced budget the default is \emph{not} the best configuration, and the ablations do not demonstrate a benefit of the graph: the identity graph (A1, no neighbours) gives \ablAOneIdentityFifteen/\ablAOneIdentityThirty/\ablAOneIdentitySixty{} mph versus \ablAZeroDefaultFifteen/\ablAZeroDefaultThirty/\ablAZeroDefaultSixty{} for the default at 15/30/60 minutes.
A random graph with the same number of edges (A2) is worse than both (\ablATwoRandomSixty{} mph at 60 minutes), so wrong edges hurt while correct edges are about neutral at this budget; this is weak support for the claim that topology matters.
The only large effects are removing the node embedding (A7) and the calendar features (A8), which hurt most at 60 minutes (\ablASevenNoNodeEmbSixty{} and \ablAEightNoTimeSixty{} mph).
Edge direction (A3), attention type (A4: GAT, GATv2, fixed weights), heads (A5), layers (A6), imputation (A9), history length (A11) and the edge bias (A12) change MAE by at most a few hundredths of a mph at 15 and 30 minutes, comparable to the seed-to-seed spread of the full runs (std \maeStgatFifteenStd{}--\maeStgatSixtyStd{} mph), so no ranking among them should be read from a single seed.
Removing dropout or weight decay (A10) helped slightly: with only 12 epochs the model is under-trained rather than over-fitted, so these results say nothing about regularisation at full training.
The three-seed comparison with the no-graph GRU (Table~\ref{tab:main}) is the stronger evidence that the graph helps at 60 minutes; the ablations could not confirm it with one seed and a short schedule.
A13 (hybrid design) is the comparison of residual, gate and average rows in Table~\ref{tab:main}.


\section{Interpretability and Maps}
Fig.~\ref{fig:err} shows where STGAT improves on XGBoost-global at 60 minutes during peaks (blue: STGAT better); gains are clustered along specific freeway segments while a few sensors are slightly worse.
In the corridor space-time diagram (Fig.~\ref{fig:st}) neither model anticipates the abrupt breakdown near 10:00 on the chosen day, which illustrates the difficulty of forecasting onset.
Attention is a qualitative probe, not causal evidence: averaged over test weekdays, the share of attention on downstream-only and upstream-only neighbours is the same at the AM peak and at night (\attAmPeakDownstreamOnly\% and \attAmPeakUpstreamOnly\% vs.\ \attNightDownstreamOnly\% and \attNightUpstreamOnly\%), so we find no peak-specific shift towards upstream or downstream sensors.
\begin{figure}[t]\centering
\includegraphics[width=\columnwidth]{error_diff_h12_peak.pdf}
\caption{Per-sensor MAE(XGB-global)$-$MAE(STGAT) at 60 min, peak hours.}\label{fig:err}
\end{figure}
\begin{figure*}[t]\centering
\includegraphics[width=0.85\textwidth]{space_time.pdf}
\caption{Corridor space-time diagram on one test day: truth and 30-minute forecasts (placed at their target time).}\label{fig:st}
\end{figure*}

\section{Simulated Control}
\textbf{This is an offline simulation on a synthetic network. It demonstrates how forecasts could inform a control policy. No live or real-world traffic control is performed or claimed.}
METR-LA sensors are on freeways, so the realistic analogue is on-ramp metering.
We simulate a 12-cell corridor with a Cell Transmission Model \cite{daganzo1994ctm} (triangular fundamental diagram, a bottleneck, three on-ramps, a 12\% capacity drop in congested cells) and compare no control, a fixed-time rate tuned on a validation day, ALINEA \cite{papageorgiou1991alinea}, a literal forecast-speed variant, and a look-ahead variant that rolls the simulator 15 minutes forward using the forecast congestion index; oracle versions use true future speeds.
Demand is synthetic, shaped by the observed congestion index (METR-LA has speed only); results are means over five test weekdays.
ALINEA reduces congested cell-minutes from \ctlNoControlCongestion{} to \ctlAlineaCongestion{} but total time spent is unchanged (\ctlNoControlTts{} vs.\ \ctlAlineaTts{} veh$\cdot$h), because delay moves to the ramps (maximum queue \ctlAlineaMax{} veh).
The look-ahead variant behaves like ALINEA (\ctlLookaheadAlineaCongestion{} cell-minutes, \ctlLookaheadAlineaTts{} veh$\cdot$h); the literal forecast-speed variant over-meters (\ctlPredictiveAlineaTts{} veh$\cdot$h; oracle: \ctlOraclePredictiveAlineaTts{}), partly because real-data forecasts do not respond to the simulator's own state (open loop).
Forecasts therefore gave no measurable benefit in this simulation.

\section{Limitations and Ethics}
One city and season; speed-only data; a fixed distance-based graph; reduced tuning and training budgets; ablations on one seed; STGAT simplified for compute; test period has more outages than training; MAPE is unstable at low speeds; the control study is synthetic and open loop.
The data are aggregate loop-detector speeds, so privacy risk is low; the main ethical risk is over-claiming real-world control benefits, which we avoid.

\section{Conclusion}
Graph structure helps most at long horizons and in congestion, but modestly, and well-tuned per-segment trees remain hard to beat; combining them with a graph model (residual or gated) gives the best results here.
Future work: full-budget training and tuning, adaptive graphs, more seeds for ablations, and flow/occupancy data (PeMS).

\bibliographystyle{IEEEtran}
\bibliography{../refs}
\end{document}
